\section{Discussion and open problems}\label{sec:discussion}

\paragraph{What the answers are.}
Without a prover, one can beat sampling exactly where the rarity of $C(x)=1$
is \emph{visible}: built from independent parts, forced literals, small
cutsets, unions, thresholds of independent blocks, or sparse monotone
overlaps, possibly disguised by local rewriting. DFSB plus the specialised
modules is sound on every circuit and costs $\poly\cdot K$ on the circuits whose
ideal $K$-explanation tree is small (\cref{thm:dfsb,thm:meta}). With a prover,
the visible structure may be replaced by \emph{provable} structure: anything a
propositional proof system can certify about where the rare inputs live
(\cref{thm:certenv}). In both settings the cost has the same shape, polynomial
in the explained surprise and exponential in the unexplained surprise, and the
exponential part is unavoidable for black-box residuals (\cref{thm:surprise}).

\paragraph{Relation to heuristic arguments.}
ARC's programme \cite{christiano2022presumption,hilton2024surprise} asks for
estimators that use a model's internals to beat sampling, and measures
explanations by surprise accounting. The present results isolate the part of
that programme that can be made \emph{sound against adversarial circuits}:
\begin{itemize}
\item the presumption of independence is certified where independence is
syntactic (variable-disjoint parts, \textsc{prod}) or where correlation
inequalities control the error (\cref{thm:janson});
\item surprise accounting becomes a cost formula for a sound verifier,
$\poly(|\pi|)+2^{\sigma_u(\pi)}$ up to polynomial factors, and both halves have
black-box lower bounds: a bit of argument saves at most a factor $2$ in
evaluations for lower bounds, and the $2^{\sigma_u}$ term is unavoidable for
two-sided intervals once the envelope is fixed (\cref{thm:surprise});
\item what soundness forbids is equally sharp: upper bounds must come from
deduction (\cref{thm:blackbox}), and no system certifies them for all circuits
(\cref{thm:ma}). A heuristic estimator that is accurate on random circuits but
not sound on adversarial ones is therefore not a counterexample to anything
here; it lives in the gap between ``no certificate exists'' and ``the estimate
is wrong''.
\end{itemize}
The rarity that defeats certification is \emph{distributed surprise}: many
weak, overlapping constraints (random CNFs below the threshold), or
pseudorandomness. Annealing and splitting estimators handle the first in
practice, and ARC-style heuristic estimators may handle both on typical
instances; neither comes with a certificate that a polynomial-time verifier can
check.

\paragraph{Open problems.}
\begin{enumerate}
\item \emph{Envelope lower bounds for distributed surprise.} Does every
polynomial-size DNNF envelope of a random 3-CNF with $1/3<\alpha<\alpha_{\rm sat}$
have $D(E)/p=2^{\Omega(n)}$? Proved here only for drop-conjunct relaxations.
\item \emph{Certified annealing.} Is there any polynomial-size, MA-checkable
certificate that a sampler for $D\mid A_1\wedge\dots\wedge A_j$ is accurate (a
checkable mixing certificate)? By \cref{thm:blackbox} it cannot be purely
statistical. This is what would make multilevel splitting sound.
\item \emph{Nested unions over approximate leaves.} Can the DNNF FPRAS be run
over DNNF skeletons whose leaves are sets given only by approximate counters,
approximate samplers and membership tests? This would give unions of unions of
products of dense blocks at any depth.
\item \emph{A natural two-sided separation.} Find a $\mathsf{UP}\cap\mathsf{coUP}$
problem not in $\BPP$ (under a standard assumption) whose unique-witness property
has polynomial-size extended Frege proofs; then the two-sided separation holds
with a standard proof checker instead of a theory like ZFC.
\item \emph{Optimal arguments.} How well can $\min_\pi\big(|\pi|+2^{\sigma_u(\pi)}\big)$
be approximated for natural certificate languages? Choosing which conjuncts to
keep is weighted set packing on supports; choosing branch variables is related to
pseudo-tree depth.
\item \emph{Networks.} Binarised or ReLU networks with rare activation patterns
are the case of most interest for alignment. Interval-bound propagation already
gives gate-wise outer approximations; can it be turned into sandwich certificates
(\cref{sec:part2}) with small looseness for trained networks?
\item \emph{Heuristic to certified.} When a presumption-of-independence estimate
is accurate, when can it be compiled into a certificate with polynomial overhead,
for instance by proving the independence it presumes, or by bounding the
correlations it ignores?
\end{enumerate}
